\begin{lemma}[Rutgers current-rung interchangeability]
Fix a positive-rung extended Tashkinov tree \(\mathcal E\) satisfying
maximum property, and assume \(\Delta(G)+1\leq k\).  Assume current
assertions (i), (iv), and (vi) uniformly for every maximum ETT with the
same rung number as \(\mathcal E\), with arbitrary certified input:
\begin{enumerate}
\item its displayed vertex set is elementary under its generating
  coloring;
\item at a terminal PE, under every proper terminal-ladder-stable
  coloring leaving exactly \(e_0\) uncolored, the terminal
  \(\gamma_n,\delta_n\)-component through the certified support meets
  the terminal ladder only at that support;
\item every such terminal-ladder-stable coloring is a
  \(\noderef{ModColoring}\) of the ETT.
\end{enumerate}
Assume also assertion (ii) for every maximum positive-rung ETT whose
rung number is strictly smaller than that of \(\mathcal E\), with
arbitrary certified input.

Write \(T_{n+1}\) for the exact final tree in the certified series of
\(\mathcal E\), \(T_n\) for its terminal ladder tree, \(D_n\) for its
terminal stored set, and \(\varphi_n\) for its terminal series coloring.
For every proper partial coloring \(\sigma\) leaving exactly \(e_0\)
uncolored which is
\((T_{n+1},D_n,\varphi_n)\)-\(\noderef{StableColoring}\), and every two
colors \(\alpha,\beta\), if
\[
 \alpha\in\overline\sigma(T_{n+1})
 \quad\hbox{or}\quad
 \beta\in\overline\sigma(T_{n+1}),
\]
then \(\alpha,\beta\) are
\(\noderef{TreeInterchangeable}\) on \(T_{n+1}\) under \(\sigma\).
\end{lemma}
\begin{proof}
Put \(F=T_{n+1}\).  We first record the symmetry and equal-color cases,
because the missing-color disjunction in the statement does not privilege
its first color.  By \(\noderef{TwoColorIncidenceGraph}\), exchanging
the names \(\alpha,\beta\) exchanges the two disjuncts in every adjacency
condition and hence gives the same simple incidence graph.  Therefore the
connected supports in \(\noderef{TwoColorComponent}\), the leaf
condition in \(\noderef{TwoColorPathComponent}\), and the quantification
in \(\noderef{TreeInterchangeable}\) are unchanged.  Thus
\[
 \noderef{TreeInterchangeable}(G,F,\sigma,\alpha,\beta)
 \quad\Longleftrightarrow\quad
 \noderef{TreeInterchangeable}(G,F,\sigma,\beta,\alpha).      \tag{1}
\]
This is a direct expansion of the three definitions, not an assumption
that the first named color is missing.

If \(\alpha=\beta\), the distinct-color clause in
\(\noderef{TwoColorPathComponent}\) is false for every root.  Hence the
universal implication in \(\noderef{TreeInterchangeable}\) is vacuous,
and the conclusion holds.  We may therefore assume
\(\alpha\ne\beta\).  Suppose for a contradiction that the desired
interchangeability fails.  Normalize the ordered pair \((a,b)\) as
follows:
\[
 (a,b)=
 \begin{cases}
  (\alpha,\beta),&
    \alpha\in\overline\sigma(F),\\
  (\beta,\alpha),&
    \beta\in\overline\sigma(F).
 \end{cases}                                                   \tag{2}
\]
In the first line use the first side of the given disjunction; use the
second line only when the first side is unavailable.  Then
\(a\ne b\), \(a\in\overline\sigma(F)\), and \(a,b\) are not
interchangeable: in the second line this last assertion is exactly the
contrapositive direction of the symmetry (1).

Apply \(\noderef{RutgersCounterexampleMinimumExists}\) to the fixed
original \(\mathcal E\), its assumed
\(\noderef{MaximumProperty}\), the uniform current assertions (i) and
(vi), the given properness and exact-one-uncolored hypotheses, the given
stability on \(F\), and the normalized colors in (2).  It supplies a
\(\noderef{RutgersClosedPresentation}\) \(\mathcal P\) relative to
the same fixed original and a \(\noderef{RutgersCounterexample}\)
\(\mathcal C\) on it such that
\[
 \noderef{RutgersCounterexampleMeasure}(\mathcal C)
 \leq
 \noderef{RutgersCounterexampleMeasure}(\mathcal D)           \tag{3}
\]
for every Rutgers presentation \(\mathcal Q\) relative to
\(\mathcal E\) and every Rutgers counterexample \(\mathcal D\) on
\(\mathcal Q\).  Notice that the minimum theorem itself constructs the
presentation and all three exact exit walks; no presentation has been
smuggled into the hypotheses of the present lemma.

Let
\(\mathcal R=\mathcal P.\mathit{corresponding}.\mathit{replacementETT}\).
The rung-agreement field of
\(\noderef{CorrespondingExtendedTashkinovTree}\) says that
\(\mathcal R\) has the same rung number as \(\mathcal E\), and
\(\mathcal P.\mathit{maximum}\) says that it has maximum property.
Apply uniform current assertion (i) to \(\mathcal R\).  Since its
displayed generating coloring is exactly
\(\mathcal P.\mathit{sigma}\), this gives
\[
 \noderef{ColoredVertexSetElementary}
       (G,\mathcal P.\mathit{sigma},V(\mathcal P.\mathit U)). \tag{4}
\]

All hypotheses of
\(\noderef{RutgersMiddleExitOutsideTerminalLadder}\) now hold: use the
degree bound, uniform current assertion (iv), the arbitrary-input
strict-lower assertion (ii), the presentation \(\mathcal P\), the
counterexample \(\mathcal C\), and (4).  If \(T_n^{\mathcal R}\) is
the actual terminal ladder tree in the dependent certified series of
\(\mathcal R\), that theorem yields
\[
                   \mathcal C.b_2\notin V(T_n^{\mathcal R}). \tag{5}
\]
The theorem handles the PE undo, the positive predecessor case, and the
rung-one base boundary internally; in particular, (5) does not appeal to
a nonexistent zero-rung ETT.

Apply \(\noderef{RutgersStrictShortening}\) to \(\mathcal E\), the
degree bound, the same uniform current assertions (i) and (vi), and the
data \(\mathcal P,\mathcal C,(5)\).  We obtain a new Rutgers closed
presentation \(\mathcal P'\) relative to the same original and a
counterexample \(\mathcal C'\) on it satisfying
\[
 \noderef{RutgersCounterexampleMeasure}(\mathcal C')
 <
 \noderef{RutgersCounterexampleMeasure}(\mathcal C).          \tag{6}
\]
Instantiate the minimum inequality (3) with exactly
\(\mathcal P',\mathcal C'\).  It gives the reverse weak inequality,
\[
 \noderef{RutgersCounterexampleMeasure}(\mathcal C)
 \leq
 \noderef{RutgersCounterexampleMeasure}(\mathcal C'),
\]
contradicting (6) in the natural-number order.

Thus the assumed failure of interchangeability is impossible for the
normalized pair.  In the first branch of (2) this is the desired ordered
pair \((\alpha,\beta)\); in the second branch, symmetry (1) returns the
same conclusion for \((\alpha,\beta)\).  Together with the equal-color
case, this proves the asserted current-rung interchangeability for every
side of the missing-color disjunction.
\end{proof}
